% ===========================================================================
%  Epidemic Network Spread -- KDD Explorations (double column) draft
%  Class: sigkddExp.cls   Bibliography: references.bib
%  Build: tectonic main.tex   (runs bibtex automatically)
%  STATUS: pre-results working draft. % TODO marks need data/figures.
% ===========================================================================
\documentclass{sigkddExp}

\usepackage{graphicx}
\usepackage{booktabs}
\usepackage{array}
\usepackage[hyphens]{url}
\usepackage{microtype}        % protrusion: tightens justification, kills most overfull boxes
\graphicspath{{figures/}}

% Let TeX stretch a little harder before declaring an overfull line.
\setlength{\emergencystretch}{2em}

% Emails are long single tokens that overflow the narrow author columns; render
% them in normal roman and let them break at . and @ so they wrap in-column.
\urlstyle{rm}

% sigkddExp.cls invokes \addauthorsection from \thebibliography but never
% defines it (used to list authors beyond the first three). We have no extra
% authors, so make it a no-op to avoid an undefined-control-sequence crash.
\providecommand{\addauthorsection}{}

\begin{document}

\title{How do Different Diseases Spread, and how does Network Science Help us Protect Populations?}

\numberofauthors{3}
\author{
\alignauthor Matilde Maria\\
       \affaddr{MSc Bioinformatics, 202509350}\\
       \email{\url{matildem.msousa@gmail.com}}
\alignauthor Lucas Caridade\\
       \affaddr{MSc Bioinformatics, 202509274}\\
       \email{\url{lucascaridade.work@gmail.com}}
\alignauthor Rafael Correia\\
       \affaddr{MSc Bioinformatics, 202207367}\\
       \email{\url{rafaelmdcorreia@gmail.com}}
}

\date{\today}
\maketitle

\begin{abstract}
How an epidemic spreads, and how best to contain it, depends jointly on the disease and on the
transport network it travels, two factors that are usually studied apart. Here we bring them
together on a realistic multimodal substrate, as a controlled and fully reproducible comparison,
delivered as a pipeline with an interactive companion app and scalable to any region. We build a
shared-node multilayer metapopulation of Europe, coupling air as diffusive relocation, land as
recurrent commuting, and water as ferries, and run five disease types parameterised from the
literature, from an immunizing measles-like SIR to a slow endemic gonorrhea-like SIS. Holding all
but the studied factor fixed, we find that targeting bridges (betweenness) consistently beats
targeting hubs (degree) across all five types, and by a wider margin as the network grows more
realistic, while grounding every flight only halves the outbreak, because land and ferry travel
keep carrying it.
\end{abstract}

\section{Introduction}

The spatial spread of infectious diseases depends both on pathogen biology and on the structure of human mobility. Characteristics such as transmissibility, infectious period, latency, mortality, reinfection, or waning immunity determine how an outbreak evolves locally, while transport infrastructures determine how infections move between populations. Air, sea, and land transport systems have increased the reach, speed, and volume of human movement, making mobility networks central to the geographical dissemination of infectious diseases \cite{tatem2006global}. In highly connected regions such as Europe, local outbreaks can therefore become regional or continental, with air travel supporting rapid long-distance seeding \cite{colizza:airline} and multiscale mobility processes, including commuting and long-range travel, shaping spatial and temporal spreading patterns \cite{balcan:multiscale}.

Mathematical epidemic models provide a controlled way to compare these mechanisms. Classical compartmental models divide the population into epidemiological states, beginning with the SIR framework of Kermack and McKendrick \cite{kermack:sir}, in which susceptible individuals become infected and are then removed or recovered. Extensions represent different disease archetypes: SEIR adds a latent exposed state, SIS represents infections without lasting immunity, SEIRS allows immunity to wane, and SEIQRD includes quarantine and disease-induced death \cite{paul2023fractional}. Such models are widely used to analyse transmission mechanisms, epidemic thresholds, and control strategies \cite{hethcote2000mathematics}. However, homogeneous compartmental models describe only the dynamics within a single population. To model spatial spread between cities, this work uses a metapopulation approach, where each city is a local population with internal epidemic dynamics and mobility flows connect cities, allowing infections to be transported across the network \cite{colizza2008metapopulation}.

Figure~\ref{fig:metapopulation_framework} summarizes this modelling logic by combining within-city epidemic dynamics with mobility flows between cities.

\begin{figure*}[tp]
    \centering
    \includegraphics[width=0.66\textwidth]{figures/metapopulation_epidemic_network_diagram.png}
    \caption[Metapopulation epidemic framework]{
    Metapopulation epidemic framework. Each circle represents a city, modelled as a local subpopulation with its own compartmental epidemic dynamics. The internal compartments illustrate a generic SEIR-type progression: $S$ denotes susceptible individuals, $E$ exposed individuals who are infected but not yet infectious, $I$ infectious individuals, and $R$ recovered or removed individuals. The dashed return arrow indicates that some models, such as SIS or SEIRS, allow individuals to become susceptible again. Arrows between cities represent mobility flows that can transport susceptible, exposed, or infected individuals and seed outbreaks elsewhere. Inspired by metapopulation epidemic models such as Colizza and Vespignani \cite{colizza2008metapopulation}.
    }
    \label{fig:metapopulation_framework}
\end{figure*}

Europe is represented here as a multilayer transport network composed of air, water, and land mobility layers. Multilayer network theory is appropriate for systems in which the same set of nodes can be connected by different types of links, each corresponding to a distinct interaction channel or transport mode \cite{kivela2014multilayer,boccaletti2014structure}. In this work, cities are nodes and transport routes are organised by layer: air links represent fast long-distance mobility, water links represent maritime and ferry-based connectivity, and land links represent regional travel and commuting. This distinction is epidemiologically relevant because mobility can operate through different mechanisms. Air and water travel are approximated as more diffusive, meaning that individuals may relocate and mix in a destination city, whereas land commuting is recurrent, meaning that individuals travel, mix temporarily, and return to their home city. Recurrent mobility can synchronize nearby cities, while long-distance diffusive mobility can seed infection in distant populations \cite{belik2011mobility,balcan:multiscale}.

Figure~\ref{fig:diffusive_recurrent_mobility} illustrates the conceptual difference between these two mobility mechanisms.

\begin{figure}[tbp]
    \centering
    \includegraphics[width=0.82\columnwidth]{figures/mobility_mechanisms_diffusive_vs_recurrent_types.png}
    \caption[Diffusive and recurrent mobility mechanisms]{
    Diffusive and recurrent mobility mechanisms. In diffusive mobility, individuals move from city $i$ to city $j$, relocate, and mix with the destination population; this approximates air travel and some water-based connections. In recurrent mobility, individuals travel from a home city $i$ to a destination $j$, mix temporarily, and return home; this represents land commuting and regional travel. Diffusive mobility can seed distant outbreaks, whereas recurrent mobility mainly couples and synchronizes nearby cities \cite{belik2011mobility,balcan:multiscale}.
    }
    \label{fig:diffusive_recurrent_mobility}
\end{figure}

Beyond mobility volume, the topology of the network influences epidemic propagation. Cities and routes are not structurally equivalent: some cities act as hubs with many connections, others as bridges between regions, and others belong to dense structural cores that may sustain spreading processes. Epidemic processes in complex networks are strongly affected by heterogeneity, connectivity, and node position \cite{pastorsatorras:review}. Moreover, geographical distance is not always the best proxy for epidemiological proximity, since strongly connected cities can be effectively close even when they are geographically distant \cite{brockmann2013hidden}.

Network science metrics can therefore support intervention design. Degree identifies the highest-traffic cities, betweenness centrality identifies the cities that bridge regions, and dense-neighbourhood measures such as triangle participation identify cities embedded in tightly connected cores that can sustain spreading \cite{kitsak:spreaders,milo:motifs}. More recent optimal-percolation methods, notably Collective Influence and non-backtracking centrality, go further by targeting the low-degree, high-impact bridge cities that raw degree overlooks \cite{morone:influence,radicchi:nonbacktracking}. These metrics can guide node interventions, interpreted as vaccination, protection, screening, or intensified surveillance in selected cities, and edge interventions, interpreted as closure or restriction of selected routes. Targeted immunization and edge-removal strategies can outperform random interventions in heterogeneous networks, although their effectiveness depends on the network structure and the timing and intensity of control \cite{pastorsatorras:immunization,bajardi2011mobility,yang2013edge}.

The objective of this work is comparative rather than predictive: the aim is not to reproduce a specific European epidemic or to provide operational forecasts, but to compare how different compartmental models, mobility layers, and intervention strategies shape simulated spreading. The central question is how the multilayer structure of European transport enables or limits different epidemic processes, and how network science metrics can guide effective containment. Combining compartmental models, a metapopulation framework, multilayer networks, and topology-based interventions, the study systematically relates human mobility, network structure, and the spatial dynamics of infectious disease spread.

\section{Methodology}
Our question is which protection strategy best contains an outbreak on the transport network
that connects Europe, and whether the answer depends on the type of disease. Answering it
fairly requires three components built in sequence: a transport network derived from open
data (\ref{building the map}), a metapopulation epidemic engine that runs on it (\ref{running outbreak}), and a set
of disease types and protection strategies evaluated on top (\ref{what to spread} and \ref{how to defend}). A
single principle constrains every component: across any comparison, only the factor under
study varies, while all other inputs are held fixed. The study is therefore comparative
rather than predictive. We do not estimate how many people a real epidemic would infect; we
quantify differences between scenarios (one strategy versus another, one disease type versus
another) under identical conditions.

\subsection{Building the map} \label{building the map}
The substrate is a weighted graph $G=(V,E)$ in which nodes $v\in V$ are cities and edges
$(i,j)\in E$ are travel routes carrying weight $w_{ij}$ proportional to traffic. We construct
$G$ with NetworkX from four open sources: air routes from OpenFlights, city populations from
GeoNames, road and rail topology from OpenStreetMap and GRIP, and ferry routes from
OpenStreetMap. Region is a parameter of construction, which is what allows us to hold the
construction \textbf{method} fixed while varying only geography.

Because a single air layer misrepresents how an epidemic actually propagates, the same node
set carries up to three layers, air, land, and water, combined into one shared-node multilayer
network (Figure 2). Edge weights are layer-specific and serve as passenger-volume proxies, but
they come from different recipes according to what data exists. Air weights are real route
frequencies from OpenFlights, and water weights are real ferry-route frequencies counted per
city pair from an OpenStreetMap ferry dump. That dump is filtered to passenger-carrying
services (OSM tag `passenger` not equal to `no`) whose endpoints are at least 5 km apart: the
first filter removes freight-only crossings, since the model concerns human movement, and the
distance threshold discards short river and canal crossings while retaining genuine straits
and sea ferries that connect distinct cities. Land has no comparably uniform relatively easy to flow data across
regions, so its weights are instead modelled with a radiation kernel
$K_{ij} = \mathrm{flux}_{ij}/\mathrm{pop}_i$ (Simini et al., 2012) computed from city
populations. The constraint we impose for fairness is not that all layers share one recipe,
but that each layer is built by a single method applied identically to every region
(OpenFlights for air, the same global ferry dump for water, the radiation kernel for land), so
that differences between regions reflect topology rather than the underlying data source;
European commuting data (Eurostat) is reserved for validating the land kernel within Europe,
not for parameterising any region. Node populations $N_i$ are taken from GeoNames where
available, with a degree-based proxy $N_i = P_0 + d(i)\,P_\mathrm{route}$ as a fallback.


\subsection{Running an outbreak on the map} \label{running outbreak}

Each city runs a compartmental model internally, while individuals move between cities along
the edges of $G$ once per day. This is a metapopulation reaction–diffusion process (Colizza
et al., 2007): a local \textbf{reaction} updates the compartments within each city, and a \textbf{mobility}
step couples cities through the network. Within city $i$ the force of infection is
$\lambda_i = \beta\, I_i / N_i$, and a daily step applies the reaction to every city and then
the mobility coupling before advancing to the next day.

The mobility step is where the multilayer map earns its keep. Treating all travel as a single
diffusive process is a known source of error (Balcan \& Vespignani, 2011), so the engine
implements two mechanisms matched to the layers:

\begin{itemize}
    \item \textbf{Diffusive mobility (air, water).} A fraction of each compartment leaves city $i$ for
  neighbour $j$ each day, proportional to the layer travel rate $\tau_\ell$ and the normalised
  edge weight $w_{ij}$, and mixes into the destination population. This is the standard
  reaction–diffusion coupling and carries a global invasion threshold in mobility rate
  (Colizza \& Vespignani, 2008).

  \item \textbf{Recurrent mobility (land, commuting).} Commuters travel to a neighbour by day, mix
  there, and return home at night, so they are not relocated. We model this by coupling the
  force of infection rather than moving population: residents of $i$ experience
  $\pi_i = \sum_j C_{ij}\, I^{*}_j / N^{*}_j$, where $C$ is the per-capita commuting matrix
  (the land kernel with a stay-home diagonal) and $I^{*}_j = \sum_i C_{ij} I_i$,
  $N^{*}_j = \sum_i C_{ij} N_i$ are the infectious and total people present at $j$. When no
  land layer is present $C$ is the identity and $\pi_i = I_i/N_i$, recovering the plain
  diffusive model.
\end{itemize}

The two mechanisms are distinct in principle: diffusive travel relocates individuals between cities, whereas recurrent commuting returns them to their home city after each day. Recurrent coupling is therefore the more appropriate representation of daily commuting behaviour.

To assess the impact of this modelling choice, we compared both mechanisms on the European air+land network (SIR with $\beta=0.32$, $\gamma=0.12$; air travel rate 0.0002; land commuting rate 0.3; seed size 2500; 210 days; averaged over seeds 0--2). The resulting epidemic trajectories were very similar, with peak active infection of approximately 71 million under recurrent mobility and 73 million under diffusive mobility. This similarity arises because the air layer already synchronises much of the network. Larger differences appear when the air layer is removed or on controlled chain networks, where diffusion produces earlier and higher peaks, although the effect remains below approximately 20%.

We therefore adopt recurrent mobility on modelling grounds rather than because it substantially changes outcomes in the baseline scenario. As a validation check, both approaches produce peak prevalence slightly below the theoretical peak expected for a fully synchronised SIR epidemic at this reproduction number (approximately 26\% of the population), indicating that the engine behaves within expectations.

Additionally, an optional extension allows transmission to occur during travel by applying the reaction dynamics to travellers for the duration of a journey. This effect is expected to be most relevant for long ferry crossings.


\subsection{Choosing what to spread} \label{what to spread}
The engine runs whichever compartmental model it is given, which makes the choice of diseases
the substantive decision. Rather than commit to named pathogens, whose exact parameters are
endlessly contestable, we test five disease \textbf{types}, each defined by its dynamical structure
and parameterised from one representative exemplar using values drawn entirely from the
literature (Table 1). The five types span the qualitative behaviours relevant to the question:
an immunizing acute epidemic (SIR), an infection with a latent period (SEIR), a lethal
infection controllable by isolation (SEIQR with a fatal compartment), an endemic infection
conferring no immunity (SIS), and a recurrent infection with waning immunity (SEIRS). Each
type adds one mechanism relative to the previous, so a difference in outcome is attributable
to that mechanism rather than to several simultaneous changes.

**Table 1: Disease types, models, and parameters.**

INSERT GO TO APPENDIX X TO SEE SOURCES

\subsection{Choosing how to defend} \label{how to defend}
With the disease types fixed, the remaining variable is the intervention, and the network
structure of Section 3.1 defines the available choices, because each intervention is a
decision about which structural elements to remove. Under a fixed budget we compare two modes.
The first is node targeting (vaccination): a fixed number of cities are rendered immune so
that they no longer transmit. The targeting rule is the object of study, and we compare no
intervention (control), random selection, selection by degree (the highest-traffic cities),
selection by betweenness (the structural bridges), and a structural rule based on $k$-core and
motif membership rather than raw connectivity (Kitsak et al., 2010). We additionally vary
coverage between low and high adherence. The second mode is edge targeting (interdiction): we
close routes rather than immunise cities, in four scenarios: (A) no intervention; (B) close
all air routes while retaining land and water; (C) close all air routes within an air-only
model; and (D) close only the top-$k$ routes ranked by degree versus by betweenness.
Scenarios B and C exist specifically to contrast the conclusion a single-layer model reaches
with the behaviour of the full multilayer network.

\subsection{Implementation and reproducibility} \label{implementation and reproducibility}
The study was implemented in Python using NumPy and NetworkX. The framework comprises three modules: data retrieval, network generation, and epidemic evaluation. Simulations are fully determined by their configuration and random seed, ensuring reproducibility across runs.

Most experiments were conducted through an interactive web-based simulator developed as part of the project. The simulator allows users to configure outbreak scenarios and observe epidemic progression through live epidemic curves and geographic visualisations, making it suitable both for experimentation and demonstration. A batch execution mode was also retained for running larger parameter sweeps and repeated trials.

Both interfaces use the same epidemic engine, configuration system, and output pipeline, ensuring that results are consistent regardless of how simulations are launched. The overall system architecture is shown in Figure x.


\section{Results}\label{sec:results}
We read every outcome in one currency, the peak of active infection, and walk the staged protocol
of Section~\ref{sec:params} one factor at a time: how each disease \emph{spreads} with no
intervention, which node-targeting rule best \emph{defends}, how that winner scales with
\emph{dose}, and what closing routes (\emph{interdiction}) achieves. All runs are on Europe, with
the disease type and the realism ladder (air $\rightarrow$ $+$water $\rightarrow$ $+$land) as the
two cross-cutting axes. Each configuration is a single deterministic run, so absolute counts are
comparative rather than predictive, and we read the differences between scenarios.

\subsection{How each disease spreads}\label{sec:res-spread}
With no intervention every one of the five types takes off on every rung of the ladder
(Figure~\ref{fig:spread}). We report the peak as a fraction of the population, so that substrates
carrying different numbers of people are compared fairly. The peak then orders cleanly by dynamical
class: the immunizing measles type (SIR) is highest, at $\approx$$61\%$ of the population on air
alone, followed by the endemic gonorrhea type (SIS, $\approx$$28\%$) and the latent COVID type
(SEIR, $\approx$$12\%$), with the waning-immunity influenza type (SEIRS) and the
isolation-controlled Ebola type (SEIQRD) lowest at $\approx$$4\%$ each. Latency delays a peak and
isolation caps it, so the two types that add those mechanisms surge least, even though every type
clears its invasion threshold~\cite{colizza:invasion}.

Climbing the realism ladder barely moves that fraction: for every disease it stays within about
three points across air, $+$water, and $+$land (measles $61\%$, $58\%$, $62\%$), which is exactly
why we normalise. In absolute counts the peak roughly doubles from air to the full multilayer
(measles $121$M $\rightarrow$ $213$M; gonorrhea $56$M $\rightarrow$ $97$M), but that rise is almost
entirely a population effect: the land layer pulls inland, non-airport cities into the network,
growing the modelled population from $199$M to $342$M. Per capita, the ladder adds people who can be
infected, not severity. Ferries even lower the peak slightly (measles $61\% \rightarrow 58\%$),
because the coastal cities they add crest out of phase and desynchronise the surge, whereas land's
role is dynamical, synchronising a larger connected population through recurrent
commuting~\cite{balcan:recurrent}. The \emph{temporal} signature, by contrast, belongs to the
disease and not the substrate: the acute measles type peaks within about three weeks (day
$17$--$21$) and burns out, while the endemic gonorrhea type is still climbing toward its plateau at
the four-year horizon.

\begin{figure}[t]
  \centering
  \includegraphics[width=\columnwidth]{F-spread.pdf}
  \caption{Uncontrolled spread. Peak active infection, as a fraction of the population, by disease
  type and substrate, up the realism ladder (air, $+$water, $+$land). Normalising by population
  removes the confound that the land layer also adds inland cities. The peak orders by dynamical
  class (immunizing $>$ endemic $>$ latent $>$ waning $\approx$ isolation-controlled) and is nearly
  flat across the ladder: realism adds people, not per-capita severity
  (Section~\ref{sec:res-spread}).}
  \label{fig:spread}
\end{figure}

\subsection{Which vaccination wins}\label{sec:res-defend}
Holding the budget at $15$ immunised cities ($75\%$ coverage) on the multilayer flagship, targeting
the structural \emph{bridges} (betweenness) yields the largest peak reduction for every one of the
five diseases (Figure~\ref{fig:defense}): measles $13.1\%$, the isolation, waning, and latent types
$12.3$--$12.7\%$, and the endemic gonorrhea type $10.8\%$. Targeting the busiest \emph{hubs}
(degree) and the dense-cluster \emph{subgraph} rule (Section~\ref{sec:defend}) tie well below it
($7$--$9.3\%$), while random immunisation is barely distinguishable from doing nothing
($0.5$--$0.8\%$). The ordering bridges~$>$~hubs~$\approx$~subgraph~$\gg$~random is uniform: it holds
for all five disease types and, in the underlying table, at all three rungs, so the win is not an
artefact of one disease or one substrate~\cite{kitsak:spreaders}.

If anything, the realism ladder sharpens this advantage rather than diluting it. The relative gap
between the peak left by hub-targeting and by bridge-targeting grows from $4$--$7.5\%$ on the
air-only network to $\approx$$12\%$ on the full multilayer for every disease (measles $7.5\%
\rightarrow 12.5\%$). Adding land and ferry layers routes yet more traffic through the same
inter-community conduits that betweenness identifies, so removing a bridge airport becomes
\emph{more} valuable, not less.

\begin{figure*}[t]
  \centering
  \includegraphics[width=0.82\textwidth]{F-defense.pdf}
  \caption{Node targeting. Peak reduction versus no vaccination for each strategy at a fixed budget
  ($15$ cities, $75\%$ coverage) on the multilayer flagship. Betweenness (bridges) wins for every
  disease type; degree (hubs) and the subgraph rule trail; random is near zero
  (Section~\ref{sec:res-defend}).}
  \label{fig:defense}
\end{figure*}

\subsection{How much budget it takes}\label{sec:res-dose}
Sweeping the winning betweenness budget from $5$ to $200$ cities on the flagship traces a
diminishing-returns curve for every disease (Figure~\ref{fig:dose}): peak reduction climbs from
$\approx$$8\%$ at $5$ cities to $30$--$34\%$ at $200$, passes $\approx$$12$--$13\%$ at the $15$-city
operating point used above, and flattens beyond $\sim$$60$--$120$ cities. Even the best rule, then,
needs a substantial budget to make a large dent, and no small set of cities contains these
outbreaks. The disease ordering is consistent at every budget: the types that confer immunity
(SIR, SEIR, SEIQRD, SEIRS) respond slightly better than the endemic SIS (at $200$ cities, Ebola
$33.8\%$ down to gonorrhea $30.1\%$), because a recovered compartment locks in the vaccination's
gain whereas SIS reinfection steadily erodes it.

\begin{figure}[t]
  \centering
  \includegraphics[width=\columnwidth]{F-dose.pdf}
  \caption{Dose--response. Peak reduction versus the number of cities vaccinated under the winning
  betweenness rule, on the flagship. Returns diminish beyond $\sim$$60$--$120$ cities; immunising
  diseases gain slightly more than the reinfecting endemic type (Section~\ref{sec:res-dose}).}
  \label{fig:dose}
\end{figure}

\subsection{Closing routes instead of cities}\label{sec:res-interdict}
The second lever closes \emph{routes} rather than immunising cities, and it is where a single-layer
model and the full multilayer part company most sharply (Figure~\ref{fig:interdiction}). Closing
\emph{all} air routes in an air-only model collapses every outbreak to $0.1$--$0.2\%$ of its full
peak (measles $0.18\%$, gonorrhea $0.12\%$): grounding flights ends the epidemic when air is the
only carrier. On the true multilayer substrate the same closure still leaves $34$--$44\%$ of the
peak standing (measles $44\%$, COVID $35\%$, Ebola $34\%$), because land commuting and ferries keep
moving people. A single-layer picture therefore overstates the benefit of a flight ban by two to
three orders of magnitude.

Closing only the top-$10$ hub airports (ranked by betweenness) is weaker still: it leaves
$56$--$70\%$ of the peak in the air-only model and $70$--$99\%$ on the multilayer, and for the slow
endemic type essentially all of it ($99\%$). Removing a handful of bridge airports barely perturbs
an outbreak that also travels by road and sea. Either way the message is the same, that the
realistic multimodal substrate flips the policy conclusion a single-layer model would reach: once
ground and ferry travel are in the picture, no air-route closure contains the epidemic on its own.

\begin{figure*}[t]
  \centering
  \includegraphics[width=0.82\textwidth]{F-interdiction.pdf}
  \caption{Edge targeting. Peak remaining (as a percentage of each substrate's own full-network
  peak) after closing all air routes (left) or the top-$10$ hub airports (right), for the air-only
  model versus the multimodal network. Closing air ends the outbreak only in the air-only world
  (its bars sit near zero); on the multimodal substrate $34$--$44\%$ of the peak survives
  (Section~\ref{sec:res-interdict}).}
  \label{fig:interdiction}
\end{figure*}


\section{Discussion}\label{sec:discussion}
The results answer the paper's question in a way that depends on how realistic the substrate is,
and they do so consistently across all five disease types. Three threads run through them: how the
disease type shapes the epidemic, why targeting bridges beats targeting hubs, and why the choice of
substrate changes what containment is even possible.

\subsection{Disease type shapes the curve}
The five types separate cleanly by their dynamical structure rather than by any single rate. The
immunizing measles type produces the tall, right-skewed, single-peaked curve characteristic of
SIR: infectious cities find abundant susceptible neighbours, incidence grows, and then the
susceptible pool is depleted irreversibly, so the outbreak burns out~\cite{kermack:sir}. The
endemic gonorrhea type has no recovered compartment, so reinfection continually replenishes
susceptibles and incidence settles at an endemic equilibrium rather than declining to zero. The
three intermediate types interpolate as expected: latency (SEIR) delays and lowers the peak, and
both isolation (SEIQRD) and short infectious windows blunt it.

One coincidence deserves care. The waning-immunity influenza type and the isolation-controlled
Ebola type land at almost the same peak height ($\sim$$8$M on air), but for mechanistic rather than
biological reasons: influenza's basic reproduction number is low ($R_0\approx1.5$), while Ebola's
is higher ($R_0\approx2$) but its isolation shortens the effective infectious period
($1/(\gamma+\kappa)\approx3.3$~days), so the peak, which depends on the infectious period as well
as $R_0$, settles at a similar level. This tells us the model is internally coherent, not that an
Ebola outbreak would resemble a flu season. It would not: the isolation type is exactly where the
well-mixed metapopulation framing is weakest (Section~\ref{sec:limitations}), so its absolute
number is the least literal of the five, and it is run on the same footing as the others only for
the sake of comparison.

\subsection{Why bridges beat hubs}
Across every disease and every rung, immunising the highest-betweenness cities reduced the peak
more than immunising the highest-degree cities, and both beat random selection. The mechanism is
structural. Because transmission on a network follows paths, and betweenness counts how many
shortest paths run through a node, removing high-betweenness cities severs the routes that connect
otherwise weakly coupled regions~\cite{kitsak:spreaders}. Removing high-degree cities instead
deletes locally busy nodes whose traffic is often redundant, since their neighbours can still reach
each other by other short paths, so the giant component and its shortest-path structure survive
more intact. Betweenness targets the bottlenecks of reachability; degree targets a volume the
network can partly route around.

The realism ladder amplifies this. The degree--betweenness gap in the resulting peak roughly
doubles from the air-only network ($4$--$7.5\%$) to the full multilayer ($\approx$$12\%$) for every
disease. The added layers create few genuinely new bridges; instead they channel the extra
recurrent and ferry traffic through the same inter-community conduits, so a high-betweenness city
bridges several layers at once and removing it disrupts all of them~\cite{sorianopanos:multiplex}.
A single-layer model, which sees only air bridges, therefore \emph{understates} how much
bridge-targeting gains, the opposite of the intuition that redundant layers should dilute
targeting.

\subsection{Realism mostly adds people, and changes what containment means}
Climbing the ladder raised the absolute peaks, but the spread results
(Section~\ref{sec:res-spread}) show that most of that rise is a population effect: the land layer
brings inland, non-airport cities into the network, and the peak, as a fraction of the population,
stays nearly flat across rungs. This is why we report that fraction rather than the raw count,
which would let the land layer's extra population masquerade as extra severity. Ferries change
little, because they reach only a handful of coastal cities; land's real contribution is dynamical,
synchronising a larger connected population through recurrent commuting~\cite{balcan:recurrent}.

The interdiction results carry the sharpest practical implication. In an air-only model, grounding
all flights ends the outbreak; on the multimodal substrate the same closure leaves a third to a
half of the peak standing, and closing only the busiest hubs does almost nothing. Whether route
closure can contain an epidemic at all thus flips with the substrate. This echoes the broader
lesson that single-layer transport models systematically mislead about containment once ground and
maritime mobility are included~\cite{balcan:multiscale}, and it is the clearest argument in the
paper for building the multilayer substrate rather than the air network alone.

\subsection{Limitations}\label{sec:limitations}
The study is comparative, not predictive, and its absolute counts should not be read as forecasts.
Three limitations bound the interpretation. First, the well-mixed metapopulation makes transmission
proportional to the local infectious fraction, which fits casual-contact respiratory and endemic
diseases far better than a contact-limited pathogen such as Ebola: real filovirus transmission
requires direct contact with symptomatic or deceased cases and is shaped by care settings that a
constant isolation rate only crudely approximates. Its peak is therefore the least trustworthy of
the five, and our metric, peak active infection, does not even reflect the isolation type's high
case-fatality ($\mu\approx0.71$), so a modest infectious peak can still be devastating. Second,
each configuration is a single deterministic run at fixed coverage, efficacy, and seed; a
stochastic ensemble and a coverage sweep would put confidence bounds on the reductions, though the
strategy \emph{ordering} is stable enough that we do not expect it to change. Third, parameters are
justified from the literature rather than calibrated to any outbreak, and the transport weights are
proxies for passenger volume, so the results are internally consistent comparisons, not fits to
data.

These bounds are also the roadmap. The same engine runs stochastic ensembles, other regions, and
finer coverage grids without modification, and the interactive simulator turns each of those into a
configuration change rather than a code change, so the findings here, bridges over hubs, realism
amplifying that gap, and the limits of route closure, form a stable core on which a more
quantitative, region-specific study can build.


\section{Conclusion}
We compared how five disease types spread across a European multilayer transport network, and which
network-science interventions best contain them, holding every factor but the one under study fixed.
Three findings recur across all five types and are resolved against a ten-seed ensemble. First, the
disease type orders the outbreak by dynamical class rather than by any single rate, while climbing
the air $\rightarrow$ $+$water $\rightarrow$ $+$land realism ladder mostly adds people rather than
per-capita severity: the peak fraction stays nearly flat, and ferries even desynchronise it. Second,
immunising the structural bridges (betweenness) beats immunising the busiest hubs (degree) for every
disease, and it is also the most geographically even choice and the one with the widest margin once
mortality rather than peak prevalence is the yardstick. Third, the multilayer substrate changes what
containment means: it makes node targeting \emph{harder}, since land and ferry redundancy routes the
outbreak around any removed city, and it overturns the interdiction conclusion a single-layer model
would reach, where grounding flights ends an air-only epidemic but leaves a third to a half of the
multimodal one standing.

Underlying these results is the network structure itself. The transport substrates are sparse but
strongly small-world, with heavy-tailed degrees and a community organisation that the land and ferry
layers sharpen, and their added bridges surface as low-degree, high-betweenness anomalous gateways:
precisely the cities betweenness targeting protects and degree targeting misses. The study is
comparative rather than predictive, and its absolute counts are not forecasts. The same engine
extends without modification to other regions, finer coverage grids, calibrated parameters, and full
transition-level stochasticity, so the stable core established here, bridges over hubs, realism
raising the cost of containment, and the limits of route closure, is a foundation on which a more
quantitative, region-specific study can build.


\bibliographystyle{abbrv}
\bibliography{references}

\appendix
\section{Disease parameters}
\label{app:sources}

Every parameter value is taken from the literature, not tuned. Each clinical rate is the
reciprocal of a measured duration of the exemplar disease ($\mathrm{rate}=1/\mathrm{duration}$);
each basic reproduction number $R_0$ comes from a separate epidemiological study; and the
transmission rate $\beta$ is derived from both, $\beta=R_0\gamma$ (or $\beta=R_0(\gamma+\kappa)$
for the isolation type). Table~\ref{tab:params-values} lists the value of every rate per type, and
Table~\ref{tab:sources} gives the source for each. Notably, measles and Ebola have similar
infectious periods ($\sim$8--10~d) yet $R_0$ differs roughly sevenfold ($\approx$15 vs
$\approx$2): the infectious period does not set $R_0$; together with $\beta$ it produces it.

\begin{table*}[t]
\centering
\footnotesize
\caption{Parameter values per disease type (rates per day; $\mathrm{rate}=1/\mathrm{duration}$,
with the duration in parentheses). No $\langle k^2\rangle/\langle k\rangle$ correction is applied,
as nodes are well-mixed cities rather than individuals. The endemic type's $\gamma$ and $R_0$ are
illustrative: with no immunity and a long, variable untreated duration, the regime matters, not
the exact number. For the lethal type $\gamma_Q\approx\gamma\approx0.10$.}
\label{tab:params-values}
\begin{tabular}{@{}ll cccccc@{}}
\toprule
\# & Type $\cdot$ model & $\gamma$ & $\sigma$ & $\kappa$ & $\omega$ & $\mu$ & $R_0$ \\
\midrule
1 & Immunizing $\cdot$ SIR    & 0.13 ($\sim$8\,d)   & ---                 & 0                 & 0                          & $\approx$0  & 12--18 \\
2 & Latent $\cdot$ SEIR       & 0.13 ($\sim$7--8\,d) & 0.20 ($\sim$5.1\,d) & 0                 & 0                          & $\sim$0.01  & 2--3.3 \\
3 & Lethal $\cdot$ SEIQRD     & 0.10 ($\sim$10\,d)  & 0.09 ($\sim$11.4\,d) & 0.20 ($\sim$5\,d) & 0                          & 0.71        & 1.5--2.5 \\
4 & Endemic $\cdot$ SIS       & $\sim$0.01--0.05    & ---                 & 0                 & $\infty$                   & 0           & low ($\gtrsim$1) \\
5 & Recurrent $\cdot$ SEIRS   & 0.21 ($\sim$4.8\,d) & 0.70 ($\sim$1.4\,d) & 0                 & $\sim$0.001 (2--7\,yr)     & $\sim$0.001 & 1.3--1.8 \\
\bottomrule
\end{tabular}
\end{table*}

\begin{table*}[t]
\centering
\footnotesize
\caption{Source for every disease-type parameter. $\gamma$, $\sigma$, $\kappa$ come from a measured
clinical duration of the exemplar; $R_0$ from a separate epidemiological study; $\beta$ is derived
from both.}
\label{tab:sources}
\begin{tabular}{@{}llll>{\raggedright\arraybackslash}p{7.4cm}@{}}
\toprule
Disease (model) & Value & Measured quantity & Source (paraphrased) \\
\midrule
Measles (SIR) & $\gamma=0.13$ & infectious $\sim$8~d & contagious 4~d before to 4~d after rash~\cite{cdc:measles} \\
              & $R_0=12$--$18$ & $R_0$ & ``often cited to be 12--18'' (review of 58 estimates)~\cite{guerra:measles} \\
              & $\omega=0$ & immunity & natural infection confers life-long immunity~\cite{cdc:measles} \\
\addlinespace
COVID (SEIR) & $\sigma=0.20$ & incubation 5.1~d & ``mean incubation period 5.1 days (95\% CI 4.5--5.8)''~\cite{lauer:incubation} \\
             & $\gamma=0.13$ & infectious $\sim$1~wk & infectiousness peaks near onset, low by $\sim$7~d~\cite{he:shedding} \\
             & $R_0=2$--$3.3$ & $R_0$ & ``pooled $R_0$ 3.32 (95\% CI 2.81--3.82)''~\cite{alimohamadi:covidr0} \\
\addlinespace
Ebola (SEIQRD) & $\sigma=0.09$ & incubation 11.4~d & ``incubation period was 11.4 days''~\cite{whoebola2014} \\
               & $\gamma,\kappa$ struct. & onset$\rightarrow$hospital $\sim$5~d & SEIHFR; community / hospital / burial transmission~\cite{legrand:ebola} \\
               & $\mu=0.71$ & case-fatality & ``case fatality rate 70.8\% (95\% CI 69--73)''~\cite{whoebola2014} \\
               & $R_0=1.5$--$2.5$ & $R_0$ (SEIR fit) & 1.51 Guinea, 2.53 Sierra Leone, 1.59 Liberia~\cite{althaus:ebolar0} \\
\addlinespace
Gonorrhea (SIS) & SIS, $\omega=\infty$ & immunity, latency & ``negligible protective immunity, negligible latent period''~\cite{hethcote:gonorrhea} \\
                & $\gamma\approx0.03$ & untreated duration & median untreated pharyngeal gonorrhea $\approx$16 weeks~\cite{barbee:gonorrhea} \\
\addlinespace
Influenza (SEIRS) & $\sigma=0.70$ & incubation 1.4~d & ``median incubation period for influenza A 1.4 days''~\cite{lessler:incubation} \\
                  & $\gamma=0.21$ & shedding 4.8~d & ``viral shedding averaged 4.80 days (95\% CI 4.31--5.29)''~\cite{carrat:influenza} \\
                  & $\omega\approx0.001$ & immunity & influenza A immunity wanes in 2--7~yr via antigenic drift~\cite{woolthuis:influenza} \\
                  & $R_0=1.3$--$1.8$ & $R_0$ (review) & median $R\approx1.80$ for the 1918 pandemic (IQR 1.47--2.27)~\cite{biggerstaff:influenzar0} \\
\bottomrule
\end{tabular}
\end{table*}

\section{Implementation architecture}
\label{app:arch}

The interactive simulator and the programmatic staged-protocol runs share one engine,
configuration system, and results pipeline (Figure~\ref{fig:arch}), so the demonstration and the
study cannot drift apart. (This figure is a working schematic and will be refined.)

\begin{figure}[t]
  \centering
  \includegraphics[width=\columnwidth]{F5-architecture.pdf}
  \caption{System architecture: one shared engine behind the interactive simulator and the
  programmatic runs, writing to a common results pipeline.}
  \label{fig:arch}
\end{figure}


\end{document}
